\documentclass[11pt, a4paper]{article}

% --- PAQUETES PRINCIPALES ---
\usepackage[utf8]{inputenc}
\usepackage[spanish, es-tabla]{babel}
\usepackage[margin=2.5cm]{geometry}
\usepackage{amsmath, amssymb, mathtools}
\usepackage{siunitx}
\usepackage{graphicx}
\usepackage{booktabs}
\usepackage{tabularx}
\usepackage[most]{tcolorbox}
\usepackage{listings}
\usepackage{xcolor}

% --- PAQUETE PARA GRÁFICAS MATEMÁTICAS ---
\usepackage{pgfplots}
\pgfplotsset{compat=1.18}

% --- CONFIGURACIÓN DE COLORES Y CÓDIGO (PYTHON) ---
\definecolor{codegreen}{rgb}{0,0.6,0}
\definecolor{codegray}{rgb}{0.5,0.5,0.5}
\definecolor{codepurple}{rgb}{0.58,0,0.82}
\definecolor{backcolour}{rgb}{0.96,0.96,0.96}

\lstdefinestyle{pythonstyle}{
    backgroundcolor=\color{backcolour},   
    commentstyle=\color{codegreen},
    keywordstyle=\color{magenta},
    numberstyle=\tiny\color{codegray},
    stringstyle=\color{codepurple},
    basicstyle=\ttfamily\footnotesize,
    breakatwhitespace=false,         
    breaklines=true,                 
    captionpos=b,                    
    keepspaces=true,                 
    numbers=left,                    
    numbersep=5pt,                  
    showspaces=false,                
    showstringspaces=false,
    showtabs=false,                  
    tabsize=4,
    language=Python
}
\lstset{style=pythonstyle}

% --- CONFIGURACIÓN DE SIUNITX ---
\sisetup{
  output-decimal-marker = {,},
  per-mode = symbol,
  inter-unit-product = \ensuremath{{}\cdot{}}
}

% --- ESTILO DE CABECERA Y ENCABEZADOS ---
\usepackage{fancyhdr}
\pagestyle{fancy}
\fancyhf{}
\lhead{\textbf{Circuitos Electrónicos III (IMT-221)}}
\rhead{Sesión 3: Ecuación Trascendente y Punto Q}
\cfoot{\thepage}
\renewcommand{\headrulewidth}{0.4pt}

% --- ENTORNOS PERSONALIZADOS ---
\newtcolorbox{aviso}[1][]{
  colback=red!5!white,
  colframe=red!75!black,
  fonttitle=\bfseries,
  title=#1
}

\begin{document}

\begin{center}
    {\LARGE \textbf{Apuntes de Cátedra: Sesión 3}}\\[0.3cm]
    {\large \textbf{Ecuación de Shockley, Recta de Carga y Punto de Operación (Q)}}\\
    \textit{Docente: M.Sc. Ing. Bernardo Quiroga Turdera \quad | \quad Fecha: 13 de agosto de 2026}
    \vspace{0.3cm}
    \hrule
\end{center}

\vspace{0.3cm}

% ==========================================
% SECCIÓN 1: BLINDAJE TEÓRICO
% ==========================================
\section{Blindaje Teórico: Mitos vs. Realidades (Feedback Sesión 2)}
Tras el trabajo tribal, consolidamos los siguientes conceptos fundamentales de la física del silicio:

\begin{center}
\renewcommand{\arraystretch}{1.5}
\begin{tabularx}{\textwidth}{@{} >{\bfseries}p{3.5cm} X X @{}}
\toprule
Misión Tribal & Concepto Clave Correcto & Mito Común a Corregir \\
\midrule
1. Dopaje y Transporte & 
La \textbf{Deriva} requiere un campo $\mathbf{E}$ externo. La \textbf{Difusión} ocurre por un gradiente de concentración ($dn/dx$). & 
\textcolor{red}{Mito:} El Silicio Tipo N tiene carga neta negativa. \textit{Realidad:} Es eléctricamente neutro; solo tiene portadores mayoritarios libres.\\

2. Unión P-N y $V_0$ & 
El Potencial de Contacto ($V_0$) nace por los iones fijos ($N_D^+, N_A^-$) en la zona de deplexión tras la recombinación. & 
\textcolor{red}{Mito:} $V_0$ se puede medir con un multímetro. \textit{Realidad:} Los potenciales de contacto de las puntas de prueba lo cancelan en equilibrio.\\

3. Shockley y Polarización & 
En polarización directa, la barrera se reduce a ($V_0 - V_D$), permitiendo que la difusión crezca exponencialmente. & 
\textcolor{red}{Mito:} $I_S$ es constante. \textit{Realidad:} Es críticamente sensible; se duplica aprox. cada \qty{10}{\celsius}.\\

4. Resistencia Dinámica & 
La resistencia de pequeña señal $r_d = \frac{\eta V_T}{I_D}$ depende inversamente de la corriente en DC. & 
\textcolor{red}{Mito:} Confundir la resistencia estática ($R_{DC} = V_D/I_D$) con la dinámica ($r_d = dV_D/dI_D$).\\
\bottomrule
\end{tabularx}
\end{center}

% ==========================================
% SECCIÓN 2: TEORÍA Y CONFLICTO NO LINEAL
% ==========================================
\section{El Conflicto No Lineal en un Circuito DC}
Dado un circuito simple con una fuente $V_{CC}$, una resistencia limitadora $R_S$ y un diodo en serie, tenemos un sistema de dos ecuaciones y dos incógnitas ($I_D, V_D$):
\begin{align}
    \text{Ecuación de la red (LVK):} \quad & I_D = \frac{V_{CC} - V_D}{R_S} \label{eq:lvk} \\
    \text{Ecuación del semiconductor:} \quad & I_D = I_S \left( \exp\left(\frac{V_D}{\eta V_T}\right) - 1 \right) \label{eq:shockley}
\end{align}
Al igualar \eqref{eq:lvk} y \eqref{eq:shockley}, obtenemos una \textbf{ecuación trascendente} que no posee una solución algebraica explícita mediante funciones elementales. 

\subsection{Método Gráfico: La Recta de Carga}
Representa todas las combinaciones posibles impuestas por la fuente y la resistencia.
\begin{itemize}
    \item \textbf{Corte en Eje $V_D$ ($I_D = 0$):} $V_{D(\max)} = V_{CC}$ (Circuito abierto).
    \item \textbf{Corte en Eje $I_D$ ($V_D = 0$):} $I_{D(\max)} = \frac{V_{CC}}{R_S}$ (Cortocircuito).
\end{itemize}

% GRÁFICA 1: RECTA DE CARGA Y PUNTO Q
\begin{figure}[h!]
    \centering
    \begin{tikzpicture}
        \begin{axis}[
            width=12cm, height=7cm,
            axis lines=middle,
            xlabel={$V_D$ (\unit{\volt})},
            ylabel={$I_D$ (\unit{\milli\ampere})},
            xmin=0, xmax=5.5,
            ymin=0, ymax=26,
            grid=both,
            grid style={dashed, gray!30},
            legend style={at={(0.5,0.95)}, anchor=north}
        ]
        
        % Recta de Carga (I_D en mA = (5 - V_D)/0.22) = 22.727 - 4.545*V_D
        \addplot[thick, red, domain=0:5] {22.727 - 4.545*x};
        \addlegendentry{Recta de Carga DC}

        % Curva Shockley (I_s = 2.52e-6 mA, eta*Vt = 0.0452375 V)
        \addplot[thick, blue, domain=0:0.74, samples=150] {2.52e-6 * (exp(x/0.0452375) - 1)};
        \addlegendentry{Ecuación de Shockley}

        % Punto Q (0.718, 19.46)
        \addplot[only marks, mark=*, black, mark size=2.5pt] coordinates {(0.718, 19.46)};
        \node[anchor=south west, font=\bfseries] at (0.718, 19.46) {$Q \, (0.718\text{ V}, 19.46\text{ mA})$};
        
        % Líneas de proyección
        \draw[dashed, black!60, thick] (0.718, 0) -- (0.718, 19.46) -- (0, 19.46);
        
        % Anotaciones de cortes
        \node[anchor=south west] at (5, 0) {$V_{CC} = 5\text{ V}$};
        \node[anchor=south west] at (0, 22.727) {$\frac{V_{CC}}{R_S}$};
        \end{axis}
    \end{tikzpicture}
    \caption{Intersección gráfica de la Ecuación Lineal (Recta de Carga) y la Ecuación No Lineal (Shockley).}
\end{figure}

% ==========================================
% SECCIÓN 3: PROBLEMA DE INGENIERÍA (NEWTON-RAPHSON)
% ==========================================
\section{Resolución del Problema de Ingeniería: Método Iterativo}

\textbf{Enunciado:} Para el circuito de protección del Hito 1 del PFI, se alimenta un diodo 1N4148 con $V_{CC} = \qty{5.0}{\volt}$ a través de $R_S = \qty{220}{\ohm}$. A temperatura ambiente (\qty{27}{\celsius} $\implies V_T = \qty{25.85}{\milli\volt}$), los parámetros son $I_S = \qty{2.52}{\nano\ampere}$ y $\eta = 1.75$. Calcule el punto $Q$ y $r_d$ iterativamente.

\vspace{0.3cm}
\noindent \textbf{Paso 1: Función Objetivo de Newton-Raphson}\\
Definimos $f(V_D) = 0$ y su derivada evaluando $\eta V_T = 1.75 \times \qty{0.02585}{\volt} = \qty{0.0452375}{\volt}$:
\begin{align*}
    f(V_D) &= I_S \left( e^{\frac{V_D}{\eta V_T}} - 1 \right) - \frac{V_{CC} - V_D}{R_S} \\
    f'(V_D) &= \frac{I_S}{\eta V_T} e^{\frac{V_D}{\eta V_T}} + \frac{1}{R_S}
\end{align*}
La actualización es: $V_D^{(k+1)} = V_D^{(k)} - \frac{f(V_D^{(k)})}{f'(V_D^{(k)})}$.

\vspace{0.3cm}
\noindent \textbf{Paso 2: Iteraciones Manuales (Semilla $V_D^{(0)} = \qty{0.7000}{\volt}$)}
\begin{align*}
    \text{\textbf{Iteración 0:} } \quad f(0.7) &= \qty{13.064}{\milli\ampere} - \qty{19.545}{\milli\ampere} = \qty{-6.481}{\milli\ampere} \\
    f'(0.7) &= \qty{0.28877}{\siemens} + \qty{0.00454}{\siemens} = \qty{0.29331}{\siemens} \\
    V_D^{(1)} &= 0.7000 - \frac{-0.006481}{0.29331} = \mathbf{\qty{0.7221}{\volt}} \\[10pt]
    \text{\textbf{Iteración 1:} } \quad f(0.7221) &= \qty{21.312}{\milli\ampere} - \qty{19.445}{\milli\ampere} = \qty{+1.867}{\milli\ampere} \\
    f'(0.7221) &= \qty{0.47111}{\siemens} + \qty{0.00454}{\siemens} = \qty{0.47565}{\siemens} \\
    V_D^{(2)} &= 0.7221 - \frac{0.001867}{0.47565} = \mathbf{\qty{0.71818}{\volt}} \\[10pt]
    \text{\textbf{Iteración 2:} } \quad f(0.71818) &= \qty{19.522}{\milli\ampere} - \qty{19.462}{\milli\ampere} = \mathbf{\qty{+0.060}{\milli\ampere}} \approx 0
\end{align*}
El valor converge a: $V_{DQ} \approx \mathbf{\qty{0.718}{\volt}}$ e $I_{DQ} = \frac{5.0 - 0.718}{220} = \mathbf{\qty{19.46}{\milli\ampere}}$.

\vspace{0.3cm}
\noindent \textbf{Paso 3: Cálculo de la Resistencia Dinámica ($r_d$)}
\begin{equation}
    r_d = \frac{\eta V_T}{I_{DQ}} = \frac{\qty{45.2375}{\milli\volt}}{\qty{19.46}{\milli\ampere}} = \mathbf{\qty{2.324}{\ohm}}
\end{equation}

% GRÁFICA 2: ZOOM A LA TANGENTE (r_d)
\begin{figure}[h!]
    \centering
    \begin{tikzpicture}
        \begin{axis}[
            width=10cm, height=6cm,
            axis lines=middle,
            xlabel={$V_D$ (\unit{\volt})},
            ylabel={$I_D$ (\unit{\milli\ampere})},
            xmin=0.69, xmax=0.74,
            ymin=5, ymax=26,
            grid=both,
            grid style={dashed, gray!30},
            legend style={at={(0.05,0.95)}, anchor=north west}
        ]
        
        % Curva Shockley (Zoom)
        \addplot[thick, blue, domain=0.69:0.74, samples=100] {2.52e-6 * (exp(x/0.0452375) - 1)};
        \addlegendentry{Curva del Diodo}

        % Tangente r_d (Pendiente m = 1/r_d = 1/2.324 Ohms = 0.430 A/V = 430.3 mA/V)
        % Ecuación: I_D - 19.46 = 430.3(V_D - 0.718) => I_D = 430.3*(x - 0.718) + 19.46
        \addplot[thick, magenta, dashed, domain=0.69:0.74] {430.3*(x - 0.718) + 19.46};
        \addlegendentry{Recta Tangente (AC)}

        % Punto Q
        \addplot[only marks, mark=*, black, mark size=2.5pt] coordinates {(0.718, 19.46)};
        \node[anchor=east, font=\bfseries] at (0.717, 19.46) {$Q$};
        
        % Explicación pendiente
        \node[anchor=west, text=magenta, font=\small, fill=white] at (0.722, 12) {Pendiente $g_d = \frac{d I_D}{d V_D} = \frac{1}{r_d}$};

        \end{axis}
    \end{tikzpicture}
    \caption{Ampliación del Punto $Q$. La resistencia dinámica $r_d$ es la inversa de la pendiente de la recta tangente.}
\end{figure}

\begin{aviso}[Conexión PFI (Hito 1)]
    En la red de protección del ADC, el diodo debe derivar picos de corriente. Su resistencia dinámica ($r_d \approx \qty{2.3}{\ohm}$) determina cuánto voltaje transitorio residual ``se le escapa'' al microcontrolador.
\end{aviso}

\newpage
% ==========================================
% SECCIÓN 4: SCRIPT PYTHON
% ==========================================
\section{Taller Computacional: Solución Numérica con Scipy}
Este \textit{script} resuelve numéricamente la ecuación trascendente mediante `fsolve`.

\begin{lstlisting}
# ==============================================================================
# UNIVERSIDAD CATOLICA BOLIVIANA "SAN PABLO" - SEDE TARIJA
# Sesion 3: Solucion Ecuacion Trascendente, Recta de Carga y Punto Q
# ==============================================================================
import numpy as np
import matplotlib.pyplot as plt
from scipy.optimize import fsolve

# 1. PARAMETROS DEL CIRCUITO Y DEL DIODO (1N4148)
Vcc = 5.0             # Voltaje fuente DC (V)
Rs = 220.0            # Resistencia limitadora (Ohms)
q = 1.602e-19         
kB = 1.3806e-23       
T_K = 27 + 273.15     
Vt = (kB * T_K) / q   # Voltaje termico (approx 25.85 mV)
Is = 2.52e-9          
eta = 1.75            

# 2. SOLUCION NUMERICA DEL PUNTO Q (fsolve)
def ecuacion_circuito(Vd):
    id_shockley = Is * (np.exp(Vd / (eta * Vt)) - 1)
    id_lvk = (Vcc - Vd) / Rs
    return id_shockley - id_lvk

Vd_guess = 0.7
Vdq = fsolve(ecuacion_circuito, Vd_guess)[0]
Idq = (Vcc - Vdq) / Rs  
rd = (eta * Vt) / Idq

print(f"Punto Q: Vdq={Vdq:.4f}V, Idq={Idq*1e3:.3f}mA, rd={rd:.3f} Ohms")
\end{lstlisting}

\end{document}